% ========== 参考文献筛选与核验流程图 - LaTeX TikZ代码版 ==========
% 可直接复制到LaTeX文档中使用

\begin{figure}[H]
\centering
\begin{tikzpicture}[
    node distance=1.8cm,
    auto,
    start/.style={
        ellipse,
        draw,
        thick,
        text width=3.5cm,
        align=center,
        minimum height=1cm,
        font=\small,
        fill=green!20
    },
    process/.style={
        rectangle,
        draw,
        thick,
        text width=4.5cm,
        align=center,
        minimum height=1.2cm,
        font=\small,
        fill=blue!10
    },
    decision/.style={
        diamond,
        draw,
        thick,
        text width=3.5cm,
        align=center,
        minimum height=1.5cm,
        font=\footnotesize,
        fill=yellow!20
    },
    exclude/.style={
        ellipse,
        draw,
        thick,
        text width=3cm,
        align=center,
        minimum height=1cm,
        font=\small,
        fill=red!20
    },
    line/.style={
        draw,
        thick,
        ->,
        >=stealth
    }
]

% 节点定义
\node [start] (start) {文献检索开始\\Literature Retrieval Initiated};

\node [process, below of=start] (search) {多数据库文献检索\\Multi-database Literature Search\\
    \scriptsize IEEE Xplore, ACM Digital Library, SpringerLink,\\
    USENIX, arXiv, Google Scholar, 官方技术文档};

\node [process, below of=search] (dedup) {去重与初步整理\\Deduplication and Initial Organization\\
    \scriptsize 删除重复文献，统一文献信息格式};

\node [process, below of=dedup] (screen) {标题与摘要相关性筛选\\Title and Abstract Relevance Screening\\
    \scriptsize 是否涉及深度学习推理、性能评估、推理系统};

\node [decision, below of=screen] (relevant) {是否与研究主题\\直接相关？\\Directly Relevant?};

\node [exclude, right of=relevant, node distance=4cm] (exclude1) {排除文献\\Excluded Records};

\node [process, below of=relevant] (verify) {来源与出版信息核验\\Source and Publication Verification\\
    \scriptsize 作者信息完整、出版来源可查、年份明确};

\node [decision, below of=verify] (reliable) {文献来源可靠且\\信息完整？\\Reliable Source?};

\node [exclude, right of=reliable, node distance=4cm] (exclude2) {排除文献\\Excluded Records};

\node [process, below of=reliable] (review) {全文审阅与方法学价值评估\\Full-text Review and Assessment\\
    \scriptsize 是否支持研究框架、具备引用价值、可复现};

\node [start, below of=review] (final) {纳入最终参考文献集合\\Included in Final Reference Set};

% 连接线
\path [line] (start) -- (search);
\path [line] (search) -- (dedup);
\path [line] (dedup) -- (screen);
\path [line] (screen) -- (relevant);
\path [line] (relevant) -- node [right] {YES} (review);
\path [line] (relevant) -- node [above] {NO} (exclude1);
\path [line] (review) -- (verify);
\path [line] (verify) -- (reliable);
\path [line] (reliable) -- node [right] {YES} (final);
\path [line] (reliable) -- node [above] {NO} (exclude2);
\path [line] (review) -- (final);

\end{tikzpicture}

\caption{参考文献筛选与核验流程\\Reference Screening and Verification Workflow}
\label{fig:ref_flow}

\end{figure}

% ========== 使用说明 ==========
% 1. 将上述代码复制到LaTeX文档中
% 2. 确保文档导入了tikz包：
%    \usepackage{tikz}
%    \usetikzlibrary{shapes,arrows,positioning}
% 3. 编译LaTeX文档
% 4. 流程图会自动生成

% ========== 特点 ==========
% - 纵向布局
% - 9个节点
% - 2个决策节点
% - 中英双语
% - 符合期刊规范
